\documentclass[11pt,a4paper]{article}

% --- Packages ---
\usepackage[margin=.75in]{geometry}
\usepackage[utf8]{inputenc}
\usepackage[T1]{fontenc}
\usepackage{lmodern} % Fix for font not found error
\usepackage{titlesec}
\usepackage{enumitem}
\usepackage{hyperref}
\usepackage{xcolor}
\usepackage{fancyhdr}
\usepackage{comment}
\usepackage{cleveref}
\usepackage[bottom]{footmisc}
\usepackage{graphicx}
\usepackage{longtable}
\usepackage{booktabs}
\usepackage{array}
\usepackage{calc}
\usepackage{etoolbox}

% Standard Pandoc setup
\providecommand{\tightlist}{%
  \setlength{\itemsep}{0pt}\setlength{\parskip}{0pt}}

\crefformat{section}{\S#2#1#3}
\crefformat{subsection}{\S#2#1#3}
\crefformat{subsubsection}{\S#2#1#3}
\crefformat{figure}{#2Figure~#1#3}
\crefformat{footnote}{#2\footnotemark[#1]#3}

\linespread{0.9} % Riduce lo spazio tra le linee

% --- Colors ---
\definecolor{primary}{RGB}{0, 0, 0}
\definecolor{links}{HTML}{0000EE}
\definecolor{blue}{HTML}{26428B}

\hypersetup{
    colorlinks=true,
    linkcolor=links,
    urlcolor=links,
    citecolor=links,
}

% --- Section Formatting ---
\titleformat{\section}
  {\normalfont\Large\bfseries\color{blue}}{\thesection}{1em}{}
\titleformat{\subsection}
{\normalfont\large\bfseries\color{blue}}{\thesubsection}{1em}{}
\titleformat{\subsubsection}
{\normalfont\bfseries\color{blue}}{\thesubsubsection}{1em}{}

% --- Header and Footer ---
\pagestyle{fancy}
\fancyhead[L]{\small Giuseppe Capasso}
\fancyhead[R]{\small devnet-workshop}
\fancyfoot[C]{\thepage}
\renewcommand{\headrulewidth}{0.4pt}

% --- Custom Title Command ---
\newcommand{\proposaltitle}[1]{
    \begin{center}
        \vspace*{0.1cm}
        {\LARGE \bfseries \color{blue} #1} \\[1.5ex]
        {\large \textbf{Giuseppe Capasso}} \\[1ex]
                \small{
            \vspace{0.1cm}
            \textbf{Materia:} devnet-workshop\\
        }
            \end{center}
    \vspace{0.3cm}
    \hrule height 0.5pt
    \vspace{0.5cm}
}

\begin{document}

\proposaltitle{A Python web-server with Docker}

\section{A Python web-server with
Docker}\label{a-python-web-server-with-docker}

\subsection{Goal}\label{goal}

We want to deploy a web server application using Docker.

\subsubsection{Integrated web-server}\label{integrated-web-server}

Did you know? Python already has an integrated web server. We can run a
web server by running the command:

\begin{verbatim}
python -m http.server 3000
\end{verbatim}

Let's put it in a Docker container. So we need a Dockerfile

\begin{verbatim}
FROM python:3.10.14-bookworm

WORKDIR /

CMD [ "python", "-m", "http.server", "3000" ]
\end{verbatim}

Build the image

\begin{verbatim}
docker build -t dtlab-simple-webserver .
\end{verbatim}

Run the container

\begin{verbatim}
docker run -p 3000:3000 dtlab-simple-webserver:latest
\end{verbatim}

\textbf{STOP and THINK}: what does \texttt{-p} flag do?

\subsection{Now it's your turn}\label{now-its-your-turn}

\subsubsection{Scenario}\label{scenario}

You are working as a DevOps engineer and you need to deploy an website.
The developer used Python and sends you a Github link. Code can be found
at this \href{https://github.com/alarmfox/dtlab-website}{link}.

Click on the \texttt{code} button and dowload the directory as a ZIP.
\includegraphics[width=5.83333in,height=2.89838in]{media/rId23.png}

You need to: * Download the source code from Github; * Create a
Dockerfile; * Running the application as a Docker container on your PC;

\paragraph{Downloading code from
Github}\label{downloading-code-from-github}

You can simply donwload someone's code from Github, by downloading a
ZIP. The code has 3 files: 1. \texttt{main.py}: a python script acting
as a webserver; 2. \texttt{index.html}: the entrypoint for every
website; 3. \texttt{requirements.txt}: a list of libraries the developer
used;

\paragraph{Creating the Dockerfile}\label{creating-the-dockerfile}

The Dockerfile is similar to the previous one, but we need to
\textbf{COPY} also the files \texttt{index.html} and
\texttt{requirements.txt}. Also we need to and a \textbf{RUN} command to
install the libraries from the \texttt{requirements.txt} file.

\paragraph{Running the application}\label{running-the-application}

Like the previous example you can the application by using the
\texttt{docker\ run} command.

\textbf{DOES IT WORK?} If not, why? Try to inspect the python script and
see if you see an error! You can fix it and re-build the image!

\end{document}
